\documentclass{homework}

\usepackage{enumitem}
\usepackage{tcolorbox}

\usepackage{tikz}
\usetikzlibrary{positioning,calc}

\usepackage{listings}

\lstset{
basicstyle=\small\ttfamily,
columns=flexible,
breaklines=true
}

\usepackage{fancyvrb}
\fvset{listparameters=\setlength{\topsep}{0pt}\setlength{\partopsep}{0pt}}

\title{03}

\begin{document}
\maketitle

\textbf{Organisational details}

\begin{itemize}
  \item Include in your report approximately how long the homework took you.
  \item You must follow the submission instructions outlined on GitLab.
  \item If the GitLab pipeline does not succeed, I will not grade your assignment (you will receive 0).
  The same applies if the Markdown report is illegibly formatted or if you miss the final submission deadline.
  \item I will provide feedback on Discord and/or GitLab.
  Check both after receiving your grade.
  \begin{itemize}
    \item My initial feedback may be quite minimal as there are so many of you.
    However, you can request additional feedback.
    Specific questions are also welcome.
    \item Do not share the feedback you receive, as this may give others an unfair advantage.
  \end{itemize}
  \item Grace clause submission deadline: 11.04.2026 23:59 EEST.
  \begin{itemize}
    \item The grace clause applies only if you submit the full work (all tasks) by the soft deadline and the GitLab pipeline succeeds.
    \item You have one week from receiving the feedback to fix and resubmit your improved work.
    You may regain up to (but not necessarily) half of the lost points.
  \end{itemize}
  \item You may share hints with each other, but not answers or code.
  Ask questions in the server if they could benefit everyone.
  \item You may -- and I encourage you to -- ask me questions.
  \item You may use AI as assistance, but not to solve the tasks themselves.
  \begin{itemize}
    \item Acceptable: e.g. checking your understanding, asking how to loop in Python.
    \item Not acceptable: e.g. generating code logic, library use, or crypto operations.
    \item If you use AI in any form, including for rewording, state in your report where and what you used it for.
    \item I recommend asking in the server or contacting me instead of using AI.
  \end{itemize}
  \item If I catch you plagiarising, copying, or using AI for code logic, I will report you to the faculty and fail you for the course.
  \begin{itemize}
    \item If during the exam you are unable to explain material for which you received homework credit, I may ask you to explain your work.
    \item If you cannot convince me that the homework is your own work, I will report you to the faculty and fail you.
  \end{itemize}
\end{itemize}

\newpage

\begin{center}
  Theory tasks
\end{center}

\textbf{Additional instructions}

\begin{itemize}
  \item \textbf{Please do not use AI for the theoretical part.}
  Challenge your understanding and look for additional written sources if the lecture and provided materials are not enough.
  Critical thinking and clear communication are skills -- use them or lose them.
  \item You must reference all external materials you rely on, including AI tools, if you choose to use them anyway.
\end{itemize}

\begin{task}{ROCA}
  Read \href{https://crocs.fi.muni.cz/public/papers/rsa_ccs17}{\emph{this blog post}} by the researchers who discovered ROCA.
  More than 750 000 Estonian ID cards were impacted by this vulnerability.

  Summarise the discovery (what, why, how to fix) in no more than five sentences.
  Then succinctly answer the following:
  \begin{itemize}
    \item Did the researchers completely break 2048-bit RSA?
    \item Are ECC keys generated on the affected chips vulnerable?
    \item How can you detect if your keys are vulnerable?
    \item What is the worst-case price for factorising a 2048-bit impacted key?
    \item Name three mitigation strategies if your hardware is impacted.
  \end{itemize}
\end{task}

\begin{task}{Bleichenbacher and PKCS\#1 v1.5}
  Read through the following article: \href{https://www.shysecurity.com/post/20160622-Bleichenbacher%20Attack}{\textit{Bleichenbacher Chosen Ciphertext Attack}}.
  Supporting materials:
  \begin{itemize}
    \item \href{https://wiki.elvis.science/index.php?title=Bleichenbacher_Attack#Steps_of_the_Attack}{\textit{Bleichenbacher Attack by the ELVIS project}}
    \item \href{https://asecuritysite.com/blog/2022-08-25_Crypto-Groundhog-Day---The-Bleichenbacher-Attack-4b85fc4accf4.html}{\textit{Cryptography Groundhog Day — The Bleichenbacher Attack}}
  \end{itemize}

  Answer the following:
  \begin{itemize}
    \item What is a padding oracle attack in cryptography?
    \item What does Bleichenbacher's attack allow an attacker to accomplish?
    \item Very generally, how does the attack work?
    \item Why does OAEP mitigate this attack? (might need a bit of additional research)
  \end{itemize}
\end{task}

\begin{task}{Opening up}
  Recall that for some secret $x\getsu\ZZ_q$ and a generator $G$, the ElGamal public key is $H\gets xG$.
  Encryption is then carried out with:
  \[
    r \getsu \ZZ_q, \ENC_H(M; r) = (rG, M + rH) = (U, V) = c
  \]
  where $r$ is the randomness which makes ElGamal a probabilistic encryption scheme.

  Suppose now that $r$ is leaked, meaning that you know the $r$ used during encryption.
  Are you then able to recover $M$ from $c$ without the secret key?
  If yes, how so?

  Can you also decrypt other messages (without gaining access to new values of $r$)?
\end{task}

\begin{task}{Following up}
  Regardless of the public-key encryption scheme, the encryption randomness must always be freshly sampled, kept secret, and discarded after encryption.

  Explain how, if the encryption randomness leaks and the set of possible messages is known (e.g. the choices \enquote{yes} or \enquote{no} in a referendum), an adversary can recover the message, regardless of the public-key encryption scheme used (e.g. RSAES-OAEP or lifted ElGamal).
\end{task}

\newpage
\setcounter{task}{0}

\begin{center}
  Practical tasks
\end{center}

\textbf{Additional instructions}

\begin{itemize}
  \item I must be able to run your program with Python 3.14 and OpenSSL v3.6.1\footnotemark{}.
  \footnotetext{You may use older versions of Python and OpenSSL as long as they are interoperable with mine.}
  \item The GitLab pipeline must succeed before you submit your assignment.
  \item You may not use third-party modules not included in \texttt{requirements.txt}.
  You are free to import any built-in module.
  If you feel the need to import another third-party module, please check with me beforehand!
  \item All errors must be handled: print your own descriptive message starting with \enquote{Handled:} and exit with error code \texttt{1}.
  Filesystem errors (e.g. file read/write errors) need not be handled.
  You \textbf{may not use} blanket try except clauses (catch all) since then I cannot know whether you anticipated the error.
  \item You must identify potential edge cases yourself.
  You may not ask an AI to implement them or identify them for you.
  \item You may change anything in the template files, unless specified otherwise, as long as the GitLab pipeline succeeds.
\end{itemize}

\begin{task}{RSA interop}
  Homework template file: \texttt{ossl\_rsa\_template.py}

  Create a Python3 program that performs RSAES-OAEP encryption/decryption by calling OpenSSL.
  Recall from the lecture that OAEP padding is required for probabilistic encryption, and to prevent some subtle attacks.
  Make sure to use OAEP with SHA-256 for padding derivation.

  Your program must provide three separate tools:
  \begin{itemize}
    \item \texttt{genpkey} -- generate an RSA private key

    Your program must generate a $3072$-bit RSA private key with OpenSSL and export it into a file in PEM format.
    Your program \emph{must not} print the key-generation progress bar to the console.

    The private key must be password protected according to the following:
    \begin{itemize}
      \item The key derivation function for obtaining the encryption key must be PBKDF2, and the parameters must be:
      \begin{itemize}
        \item \texttt{HMAC-SHA512} for the PRF,
        \item an iteration count of $210000$ (OWASP recommendation).
      \end{itemize}
      \item The key must be encrypted with AES-256 in CBC mode and use the PKCS\#8 format.
      \item At no point should the private key be written unencrypted to the disk.
    \end{itemize}

    \begin{tcolorbox}[title=Note]
      New versions of OpenSSL (> 3.0) should use PKCS\#8 by default if you export an encrypted private key in PEM, but do not allow for the fine-grained control needed for this task.

      Hint: Can you find a tool in the OpenSSL documentation which allows more fine-grained control over PKCS\#8 encrypted private keys? 
    \end{tcolorbox}

    You can see what the expected private key structure looks like by dumping the ASN.1 structure of the example private key.

    Example invocation:
    \begin{Verbatim}
ossl_rsa.py genpkey --out key.pem --pass secret
    \end{Verbatim}

    \item \texttt{pkey} -- public key derivation
    
    Your program must extract the public key from a private key and export it into a file in PEM format by using OpenSSL.

    Example invocation:
    \begin{Verbatim}
ossl_rsa.py pkey --in key.pem \
    --out pub.pem --pass secret
    \end{Verbatim}

    \item \texttt{pkeyutl} --- file encryption and decryption
    
    Your program must output the encryption or decryption of the specified input file with the specified key to the specified output file using OpenSSL.

    Example invocation for encryption:
    \begin{Verbatim}
ossl_rsa.py pkeyutl --encrypt \
    --inkey pub.pem --in message.txt --out data.enc
    \end{Verbatim}

    Example invocation for decryption:
    \begin{Verbatim}
ossl_rsa.py pkeyutl \
    --inkey key.pem --in data.enc --out data.dec \
    --pass secret
    \end{Verbatim}
  \end{itemize}

  Make sure to handle errors by printing an informational message and exiting with an error code (\texttt{sys.exit(1)}) without printing the exception trace.
  Your program must not crash if missing parameters are provided, the wrong password is provided, OpenSSL fails, etc.
  Some error handling and hints are part of the template, but not all of it.

  All printed error messages must start with the string \enquote{Handled:}.

  You do not however have to (but you may) check for file I/O errors, i.e. you may assume that if you provide a filename, then that filename does exist or can be written to.

  \begin{tcolorbox}[title=Reminder]
    In practice, providing secrets (e.g. keys, passwords, etc.) directly on the CLI as an argument to a program is a terrible idea, since they typically get logged into your terminal history!
  \end{tcolorbox}
\end{task}

\begin{task}{Nifty homomorphism}
  Task template file: \texttt{homomorphic\_template.py}

  \textit{Preface.}
  The legendary cryptographer Mr. Cry has developed the following authentication protocol:
  \begin{itemize}
    \item Authorised users know the secret authentication code.
    \item To gain access to the server, users encrypt this code (with ElGamal) to obtain the \emph{token}, submit the token to the server, and the server decrypts it.
    If the decrypted result corresponds to the authentication code, the user is granted access.
    \item To avoid authentication token re-use, the server remembers all previously used tokens.
    If you submit a previously used token, the server will deny you access.
    \item To mitigate \enquote{double-submission} (you submit the same token twice e.g. due to bad latency), the server displays the previously used token to the user.
  \end{itemize}

  Your goal is to forge a token and thus gain access.

  \textit{Requirements.}
  Implement the function
  \begin{Verbatim}
octets_to_point(octets: bytes, curve: Curve) -> Point
  \end{Verbatim}
  that parses a SEC~1 uncompressed point byte string representation\footnote{\texttt{0x04 || xCoordBytes || yCoordBytes} as specified in \href{https://www.secg.org/sec1-v2.pdf\#subsubsection.2.3.3}{SEC~1 v2.0 \S 2.3.3}.} and returns the corresponding point on the given curve.

  Implement the function
  \begin{Verbatim}
point_to_octets(p: Point) -> bytes
  \end{Verbatim}
  that serialises a point into the SEC~1 uncompressed point representation.

  The \texttt{Point} and \texttt{Curve} types may be taken from the \texttt{fastecdsa} (preferred) or \texttt{tinyec} library.
  If the library has appropriate methods for (de)serialisation, you may use those in the functions.

  Your program takes as input an EC ElGamal ciphertext (the token), consisting of two hex-encoded points ($U$, $V$), and a file containing the PEM-encoded EC public key in the standard\footnote{\href{https://datatracker.ietf.org/doc/html/rfc5480}{RFC 5480}, with \href{https://datatracker.ietf.org/doc/html/rfc3279}{RFC 3279} and \href{https://datatracker.ietf.org/doc/html/rfc5280\#section-4.1.2.7}{RFC 5280} being related.} format.
  The program outputs a fresh token, printing $U$ and $V$ as hex-encoded points on separate lines.\footnote{This is taken care of by the template.}
  E.g.
  \begin{Verbatim}
homomorphic.py -u 048be2...b6e9 -v 0462c3...4af0 -k pub.pem
04ff6d...e83a
0433f7...8aa7
  \end{Verbatim}
  without the dots in the middle.

  The program must work with the P-384 curve, and you may hardcode that curve value in the main function.
  However, if an unsupported public key is provided, your program must handle this without crashing.
  
  Your task succeeds if:
  \begin{itemize}
    \item $\DEC_x\bigl((U_\mathit{challenge}, V_\mathit{challenge})\bigr) = \DEC_x\bigl((U_\mathit{you}, V_\mathit{you})\bigr)$, and
    \item $(U_\mathit{challenge}, V_\mathit{challenge}) \neq (U_\mathit{you}, V_\mathit{you})$.
  \end{itemize}

  You can get a test token from the Discord bot and submit your answer to it to validate your solution.
\end{task}

\begin{task}{Posting the quantum}
  Task template file: \texttt{mlkem\_template.py}

  The golden rule of cryptography is that you should not roll your own cryptography.
  But what do we do when libraries and tooling do not yet implement everything we need?

  Realistically, we should wait for appropriate library and tooling support to become available, or hire experts, neither of which is an option for this homework assignment.

  Implement the \texttt{exchange} tool using OpenSSL's ML-KEM support to encapsulate a shared secret under the provided ML-KEM public key.

  Implement the \texttt{decrypt} tool that uses HKDF to derive a 256-bit AES key from the shared secret, and uses it to decrypt an AES-GCM authenticated ciphertext.
  HKDF should be instantiated with SHA-384, the context info string \texttt{ITC8280 Assignment 3}, and no salt.

  Supporting files:
  \begin{itemize}
    \item \texttt{mlkem.pub.pem} -- an example ML-KEM public key
    \item \texttt{shared.bin} -- a shared secret (in practice, this should be secret)
    \item \texttt{ciphertext.bin} -- an authenticated ciphertext encrypted with a key derived from the shared secret
  \end{itemize}
\end{task}

\end{document}
